% Restrained TikZ redraws for Lesson 32.

\newcommand{\arbitraryCapacitorDiagram}{%
\begin{center}
\begin{tikzpicture}[x=1cm,y=1cm,>=Latex]
  \fill[PaleBlue] plot[smooth cycle,tension=0.78] coordinates
    {(-3.7,0.1) (-3.3,1.35) (-2.05,1.55) (-1.45,0.55) (-1.8,-0.75) (-3.05,-1.25) (-3.95,-0.55)};
  \draw[DeepBlue,very thick] plot[smooth cycle,tension=0.78] coordinates
    {(-3.7,0.1) (-3.3,1.35) (-2.05,1.55) (-1.45,0.55) (-1.8,-0.75) (-3.05,-1.25) (-3.95,-0.55)};
  \fill[SoftGray] plot[smooth cycle,tension=0.78] coordinates
    {(3.7,0.15) (3.15,1.45) (1.85,1.25) (1.4,0.15) (1.95,-1.05) (3.35,-1.25) (4.0,-0.45)};
  \draw[Ink,very thick] plot[smooth cycle,tension=0.78] coordinates
    {(3.7,0.15) (3.15,1.45) (1.85,1.25) (1.4,0.15) (1.95,-1.05) (3.35,-1.25) (4.0,-0.45)};
  \foreach \p in {(-3.35,0.65),(-2.65,1.05),(-2.15,0.15),(-3.1,-0.55)}
    \node[DeepBlue] at \p {$+$};
  \foreach \p in {(3.3,0.7),(2.45,0.95),(1.95,0.05),(2.85,-0.55)}
    \node[Ink] at \p {$-$};
  \draw[component] (-1.65,0.82) .. controls (-0.5,1.5) and (0.5,1.5) .. (1.65,0.82);
  \draw[component] (-1.45,0.15) .. controls (-0.45,0.65) and (0.45,0.65) .. (1.45,0.15);
  \draw[component] (-1.62,-0.55) .. controls (-0.5,-1.15) and (0.5,-1.15) .. (1.72,-0.55);
  \node[DeepBlue] at (-2.7,-1.65) {$+Q,\ \varphi_+$};
  \node[Ink] at (2.7,-1.65) {$-Q,\ \varphi_-$};
  \draw[<->,TextGray] (-1.05,-1.9)--node[below] {$U=\varphi_+-\varphi_-$}(1.05,-1.9);
\end{tikzpicture}

{\small\color{TextGray}图 32.1\quad 任意形状的两块导体组成双极电容器}
\end{center}}

\newcommand{\capacitorGeometryDiagram}{%
\begin{center}
\begin{tikzpicture}[scale=0.82,>=Latex]
  % parallel plates
  \draw[DeepBlue,very thick] (-5.2,0.8)--(-2.6,0.8);
  \draw[DeepBlue,very thick] (-5.2,-0.8)--(-2.6,-0.8);
  \foreach \x in {-4.9,-4.4,...,-2.9}{
    \node[DeepBlue] at (\x,0.48) {$+$}; \node[Ink] at (\x,-0.48) {$-$};
    \draw[component] (\x,0.32)--(\x,-0.32);
  }
  \draw[<->] (-5.55,-0.8)--node[left] {$d$}(-5.55,0.8);
  \node[TextGray] at (-3.9,-1.45) {平行板};
  % spherical
  \draw[curve] (0,0) circle (0.73); \draw[curve] (0,0) circle (1.28);
  \draw[<->] (0,0)--node[above] {$R_1$}(0.73,0);
  \draw[<->] (0,-0.28)--node[below] {$R_2$}(1.28,-0.28);
  \node[TextGray] at (0,-1.45) {同心球};
  % coaxial
  \draw[curve] (3.6,0) circle (0.62); \draw[curve] (3.6,0) circle (1.20);
  \foreach \a in {0,45,...,315}{\draw[component] ({3.6+0.66*cos(\a)},{0.66*sin(\a)})--({3.6+1.12*cos(\a)},{1.12*sin(\a)});}
  \node at (3.6,0) {$a$}; \node[above right] at (4.35,0.75) {$b$};
  \node[TextGray] at (3.6,-1.45) {同轴圆柱截面};
\end{tikzpicture}

{\small\color{TextGray}图 32.2、32.3、32.10、32.11\quad 平行板、同心球与同轴圆柱电容器（合并重绘）}
\end{center}}

\newcommand{\seriesParallelDiagram}{%
\begin{center}
\begin{tikzpicture}[scale=0.86,>=Latex]
  % parallel
  \draw[thick] (-5.0,1.25)--(-4.2,1.25) (-5.0,-1.25)--(-4.2,-1.25);
  \draw[thick] (-4.2,1.25)--(-4.2,-1.25) (-1.5,1.25)--(-1.5,-1.25);
  \draw[thick] (-4.2,0.72)--(-3.3,0.72) (-2.4,0.72)--(-1.5,0.72);
  \draw[very thick,DeepBlue] (-3.3,0.38)--(-3.3,1.06) (-2.4,0.38)--(-2.4,1.06);
  \draw[thick] (-4.2,-0.72)--(-3.3,-0.72) (-2.4,-0.72)--(-1.5,-0.72);
  \draw[very thick,DeepBlue] (-3.3,-1.06)--(-3.3,-0.38) (-2.4,-1.06)--(-2.4,-0.38);
  \node at (-2.85,1.28) {$C_1$}; \node at (-2.85,-1.28) {$C_2$};
  \node[TextGray] at (-3.25,-1.85) {并联：电压相同};
  % series
  \draw[thick] (0.8,0)--(1.8,0) (2.55,0)--(3.45,0) (4.2,0)--(5.2,0);
  \draw[very thick,DeepBlue] (1.8,-0.55)--(1.8,0.55) (2.55,-0.55)--(2.55,0.55);
  \draw[very thick,DeepBlue] (3.45,-0.55)--(3.45,0.55) (4.2,-0.55)--(4.2,0.55);
  \node at (2.18,0.95) {$C_1$}; \node at (3.82,0.95) {$C_2$};
  \node[TextGray] at (3.0,-1.85) {串联：中间孤立导体净电荷不变};
\end{tikzpicture}

{\small\color{TextGray}图 32.4、32.5\quad 电容器的并联与串联（合并重绘）}
\end{center}}

\newcommand{\polarizationDiagram}{%
\begin{center}
\begin{tikzpicture}[scale=0.87,>=Latex]
  \fill[PaleBlue] (-3.3,-1.5) rectangle (3.3,1.5);
  \draw[TextGray,thick] (-3.3,-1.5) rectangle (3.3,1.5);
  \foreach \x in {-2.6,-1.3,0,1.3,2.6}{
    \foreach \y in {-0.85,0,0.85}{
      \node[Ink] at (\x-0.16,\y) {$-$}; \node[DeepBlue] at (\x+0.16,\y) {$+$};
    }
  }
  \foreach \y in {-1.05,-0.35,0.35,1.05}{
    \node[Ink] at (-3.55,\y) {$-$}; \node[DeepBlue] at (3.55,\y) {$+$};
  }
  \draw[vector] (-1.0,2.05)--(1.0,2.05) node[midway,above] {$\vect E$};
  \draw[vector] (-1.0,-2.05)--(1.0,-2.05) node[midway,below] {$\vect P$};
  \node[TextGray,align=left,anchor=west] at (4.25,0.45)
    {分子偶极矩沿外场排列\\两端出现束缚面电荷};
\end{tikzpicture}

{\small\color{TextGray}图 32.9\quad 外场作用下电介质的极化与束缚面电荷}
\end{center}}

\newcommand{\reciprocitySphereDiagram}{%
\begin{center}
\begin{tikzpicture}[x=1cm,y=1cm,>=Latex]
  \draw[curve] (-2.5,0) circle (1.18);
  \fill (-2.5,0) circle (1.8pt) node[below=4pt] {$O$};
  \fill[DeepBlue] (2.5,0) circle (2.6pt) node[above=5pt] {$P$};
  \draw[<->] (-2.5,0.25)--node[above] {$r$}(2.5,0.25);
  \node[align=center] at (-2.5,-1.65) {状态 I：球带单位电荷\\状态 II：$P$ 点放电荷 $q$};
  \node[align=left,anchor=west,TextGray] at (4.0,0.45)
    {$\displaystyle \sum_i Q_i\varphi_i'=\sum_iQ_i'\varphi_i$\\[4pt]
     交换“源”与“测量位置”，\\倒易关系保持不变};
\end{tikzpicture}

{\small\color{TextGray}图 32.6\quad Green 倒易定理：导体球与球外一点的两种静电状态}
\end{center}}

\newcommand{\reciprocityConductorsDiagram}{%
\begin{center}
\begin{tikzpicture}[x=1cm,y=1cm,>=Latex]
  \fill[PaleBlue] plot[smooth cycle,tension=0.8] coordinates
    {(-4.1,0.2) (-3.5,1.25) (-2.15,1.05) (-1.65,0) (-2.3,-1.1) (-3.7,-0.95)};
  \draw[DeepBlue,thick] plot[smooth cycle,tension=0.8] coordinates
    {(-4.1,0.2) (-3.5,1.25) (-2.15,1.05) (-1.65,0) (-2.3,-1.1) (-3.7,-0.95)};
  \fill[SoftGray] plot[smooth cycle,tension=0.8] coordinates
    {(4.05,0.15) (3.35,1.2) (2.0,0.9) (1.55,-0.15) (2.35,-1.2) (3.75,-0.85)};
  \draw[Ink,thick] plot[smooth cycle,tension=0.8] coordinates
    {(4.05,0.15) (3.35,1.2) (2.0,0.9) (1.55,-0.15) (2.35,-1.2) (3.75,-0.85)};
  \node[align=center] at (-2.85,0) {$Q_1$\\$\varphi_1$};
  \node[align=center] at (2.85,0) {$Q_2$\\$\varphi_2$};
  \draw[<->,TextGray] (-1.15,0)--node[above] {静电感应}(1.15,0);
  \node[align=center,TextGray] at (0,-1.7)
    {第二状态用 $Q_1',Q_2',\varphi_1',\varphi_2'$ 表示；导体几何保持不变};
\end{tikzpicture}

{\small\color{TextGray}图 32.7\quad Green 倒易定理：两个任意导体的两种静电平衡状态}
\end{center}}

\newcommand{\dielectricExperimentDiagram}{%
\begin{center}
\begin{tikzpicture}[x=0.86cm,y=0.86cm,>=Latex]
  \foreach \shift/\lab/\needle in {-4.3/{未插入介质：$C_0,U_0$}/32,4.3/{插入介质：$C=\varepsilon_rC_0$}/-8}{
    \begin{scope}[shift={(\shift,0)}]
      \draw[Ink,very thick] (-1.15,1.0)--(-1.15,3.1);
      \draw[Ink,very thick] (1.15,1.0)--(1.15,3.1);
      \node[DeepBlue] at (-1.45,2.1) {$+$}; \node[Ink] at (1.45,2.1) {$-$};
      \ifdim\shift pt>0pt
        \fill[DeepBlue!18] (-0.95,1.15) rectangle (0.95,2.95);
        \draw[DeepBlue,thick] (-0.95,1.15) rectangle (0.95,2.95);
        \node[DeepBlue,rotate=90] at (0,2.05) {电介质};
      \fi
      \draw[Ink,thick] (-1.15,1.0)--(-1.15,0.45)--(-0.75,0.45);
      \draw[Ink,thick] (1.15,1.0)--(1.15,0.45)--(0.75,0.45);
      \draw[rounded corners=3pt,fill=SoftGray,draw=Ink,thick] (-0.75,-1.5) rectangle (0.75,0.55);
      \draw[Ink] (0,-0.25) circle (0.48);
      \draw[DeepBlue,very thick] (0,-0.25)--({0.38*cos(\needle)},{-0.25+0.38*sin(\needle)});
      \foreach \a in {20,50,...,160}{\draw[TextGray] ({0.38*cos(\a)},{-0.25+0.38*sin(\a)})--({0.45*cos(\a)},{-0.25+0.45*sin(\a)});}
      \node[align=center] at (0,-2.05) {\lab};
    \end{scope}
  }
  \draw[vector] (-1.15,0.2)--(1.15,0.2) node[midway,above] {插入介质};
\end{tikzpicture}

{\small\color{TextGray}图 32.8\quad 保持电荷不变时，插入电介质使电势差减小、电容增大}
\end{center}}

\newcommand{\touchingSpheresDiagram}{%
\begin{center}
\begin{tikzpicture}[scale=0.92,>=Latex]
  \draw[curve] (-1.3,0) circle (1.3); \draw[curve] (1.3,0) circle (1.3);
  \fill (0,0) circle (1.4pt) node[below] {接触点};
  \draw[<->] (1.3,0)--node[above] {$R$}(2.6,0);
  \foreach \x/\r in {-1.75/2.3,-0.95/1.8,-0.48/1.25,0.48/1.25,0.95/1.8,1.75/2.3}{
    \fill[DeepBlue] (\x,0) circle (\r pt);
  }
  \node[TextGray] at (0,-2.05) {接触极限可由双球坐标或无穷像电荷序列求得};
\end{tikzpicture}

{\small\color{TextGray}图 32.12\quad 两个等半径相切球连成一个等势导体}
\end{center}}

\newcommand{\twoCylinderDiagram}{%
\begin{center}
\begin{tikzpicture}[scale=0.9,>=Latex]
  \draw[curve] (-2.2,0) circle (0.85); \draw[curve] (2.2,0) circle (0.85);
  \node at (-2.2,0) {$+\lambda$}; \node at (2.2,0) {$-\lambda$};
  \fill ( -2.2,0) circle (1.3pt); \fill (2.2,0) circle (1.3pt);
  \draw[<->] (-2.2,-1.28)--node[fill=white,inner sep=1pt] {$L$}(2.2,-1.28);
  \draw[<->] (-2.2,0)--node[above] {$R$}(-1.35,0);
  \draw[DeepBlue,thick] (-1.55,0.8) .. controls (0,2.0) .. (1.55,0.8);
  \draw[DeepBlue,thick] (-1.55,-0.8) .. controls (0,-2.0) .. (1.55,-0.8);
  \draw[DeepBlue,thick] (-1.35,0.25) .. controls (0,0.78) .. (1.35,0.25);
  \draw[DeepBlue,thick] (-1.35,-0.25) .. controls (0,-0.78) .. (1.35,-0.25);
  \node[TextGray] at (0,-2.35) {$\cosh\tau_0=L/(2R)$，两圆柱面对应 $\tau=\pm\tau_0$};
\end{tikzpicture}

{\small\color{TextGray}图 32.13\quad 两根等半径平行圆柱导体的横截面}
\end{center}}

\newcommand{\splitDielectricDiagram}{%
\begin{center}
\begin{tikzpicture}[scale=0.83,>=Latex]
  % side-by-side
  \draw[DeepBlue,very thick] (-5.1,1.0)--(-1.9,1.0) (-5.1,-1.0)--(-1.9,-1.0);
  \fill[PaleBlue] (-5.1,-0.95) rectangle (-3.5,0.95);
  \draw[guide] (-3.5,-1.0)--(-3.5,1.0);
  \node at (-4.3,0) {$\varepsilon$}; \node at (-2.7,0) {$\varepsilon_0$};
  \node[TextGray] at (-3.5,-1.55) {按面积分割：并联};
  % layered
  \draw[DeepBlue,very thick] (1.2,1.0)--(4.4,1.0) (1.2,-1.0)--(4.4,-1.0);
  \fill[PaleBlue] (1.2,0) rectangle (4.4,0.95);
  \draw[guide] (1.2,0)--(4.4,0);
  \node at (2.8,0.5) {$\varepsilon$}; \node at (2.8,-0.5) {$\varepsilon_0$};
  \node[TextGray] at (2.8,-1.55) {按厚度分割：串联};
\end{tikzpicture}

{\small\color{TextGray}图 32.14\quad 部分填充介质的两种等效方式}
\end{center}}

\newcommand{\dielectricPullDiagram}{%
\begin{center}
\begin{tikzpicture}[scale=0.88,>=Latex]
  \draw[DeepBlue,very thick] (-3.8,0.8)--(3.8,0.8) (-3.8,-0.8)--(3.8,-0.8);
  \fill[PaleBlue] (-3.8,-0.72) rectangle (0.6,0.72);
  \draw[TextGray,thick] (-3.8,-0.72) rectangle (0.6,0.72);
  \draw[<->] (-3.8,-1.22)--node[fill=white,inner sep=1pt] {$x$}(0.6,-1.22);
  \draw[<->] (0.6,-1.22)--node[fill=white,inner sep=1pt] {$a-x$}(3.8,-1.22);
  \draw[<->] (4.15,-0.8)--node[right] {$d$}(4.15,0.8);
  \draw[vector] (0.55,0)--(-1.1,0) node[midway,above] {$\vect F$};
  \node[TextGray] at (0,-1.95) {介质进入越深，等效电容越大，电场把介质继续拉入};
\end{tikzpicture}

{\small\color{TextGray}图 32.15\quad 部分插入平行板电容器的介质板}
\end{center}}

\newcommand{\dielectricSphereDiagram}{%
\begin{center}
\begin{tikzpicture}[scale=0.88,>=Latex]
  \fill[PaleBlue] (0,0) circle (1.35); \draw[curve] (0,0) circle (1.35);
  \foreach \a in {-60,-30,0,30,60}{
    \node[Ink] at ({-1.08*cos(\a)},{1.08*sin(\a)}) {$-$};
    \node[DeepBlue] at ({1.08*cos(\a)},{1.08*sin(\a)}) {$+$};
  }
  \draw[vector] (-4.3,1.95)--(-2.8,1.95) node[midway,above] {$\vect E_0$};
  \draw[DeepBlue,thick,-{Latex[length=2mm]}] (-4.2,1.72)--(4.2,1.72);
  \draw[DeepBlue,thick,-{Latex[length=2mm]}] (-4.2,-1.72)--(4.2,-1.72);
  \draw[DeepBlue,thick] (-4.2,0.92) .. controls (-2.6,0.92) and (-1.8,0.62) .. (-1.26,0.48);
  \draw[DeepBlue,thick,-{Latex[length=2mm]}] (-1.26,0.48)--(1.26,0.48);
  \draw[DeepBlue,thick,-{Latex[length=2mm]}] (1.26,0.48) .. controls (1.8,0.62) and (2.6,0.92) .. (4.2,0.92);
  \draw[DeepBlue,thick] (-4.2,-0.92) .. controls (-2.6,-0.92) and (-1.8,-0.62) .. (-1.26,-0.48);
  \draw[DeepBlue,thick,-{Latex[length=2mm]}] (-1.26,-0.48)--(1.26,-0.48);
  \draw[DeepBlue,thick,-{Latex[length=2mm]}] (1.26,-0.48) .. controls (1.8,-0.62) and (2.6,-0.92) .. (4.2,-0.92);
  \draw[DeepBlue,thick] (-4.2,0)--(-1.35,0);
  \draw[DeepBlue,thick,-{Latex[length=2mm]}] (-1.35,0)--(1.35,0);
  \draw[DeepBlue,thick,-{Latex[length=2mm]}] (1.35,0)--(4.2,0);
  \node at (0,-0.82) {$\varepsilon$};
  \node[TextGray] at (0,-2.35) {$\vect E_{\rm in}=\dfrac{3\varepsilon_0}{\varepsilon+2\varepsilon_0}\vect E_0$};
\end{tikzpicture}

{\small\color{TextGray}解图 32.6-A\quad 匀强电场中的线性电介质球}
\end{center}}

\newcommand{\chargingDiagram}{%
\begin{center}
\begin{tikzpicture}[scale=0.9,>=Latex]
  \draw[thick] (-3.8,0)--(-2.9,0);
  \draw[thick] (-2.9,-0.55)--(-2.9,0.55); \draw[very thick,DeepBlue] (-2.55,-0.35)--(-2.55,0.35);
  \node[below] at (-2.72,-0.62) {$\mathcal E$};
  \draw[thick] (-2.55,0)--(-1.7,0);
  \draw[thick,decorate,decoration={zigzag,segment length=6pt,amplitude=2.5pt}] (-1.7,0)--(0.2,0);
  \node[above] at (-0.75,0.35) {$R$};
  \draw[thick] (0.2,0)--(1.25,0);
  \draw[very thick,DeepBlue] (1.25,-0.55)--(1.25,0.55) (1.85,-0.55)--(1.85,0.55);
  \node[above] at (1.55,0.75) {$C$};
  \draw[thick] (1.85,0)--(2.8,0)--(2.8,-1.35)--(-3.8,-1.35)--(-3.8,0);
  \node[TextGray,align=left,anchor=west] at (4.0,-0.1)
    {电源功 $C\mathcal E^2$\\电场能 $C\mathcal E^2/2$\\焦耳热 $C\mathcal E^2/2$};
\end{tikzpicture}

{\small\color{TextGray}解图 32.8-A\quad 电阻参与的电容器充电过程}
\end{center}}

\newcommand{\starCapacitorDiagram}{%
\begin{center}
\begin{tikzpicture}[scale=0.94,>=Latex]
  \fill[DeepBlue] (0,0) circle (2pt) node[below right] {$\varphi_0$};
  \foreach \a/\c/\p in {90/C_1/\varphi_1,210/C_2/\varphi_2,330/C_3/\varphi_3}{
    \draw[thick] (0,0)--({0.95*cos(\a)},{0.95*sin(\a)});
    \draw[very thick,DeepBlue] ({0.78*cos(\a)-0.35*sin(\a)},{0.78*sin(\a)+0.35*cos(\a)})--({0.78*cos(\a)+0.35*sin(\a)},{0.78*sin(\a)-0.35*cos(\a)});
    \draw[very thick,DeepBlue] ({1.12*cos(\a)-0.35*sin(\a)},{1.12*sin(\a)+0.35*cos(\a)})--({1.12*cos(\a)+0.35*sin(\a)},{1.12*sin(\a)-0.35*cos(\a)});
    \draw[thick] ({1.12*cos(\a)},{1.12*sin(\a)})--({2.2*cos(\a)},{2.2*sin(\a)});
    \node at ({1.48*cos(\a)+0.45*sin(\a)},{1.48*sin(\a)-0.45*cos(\a)}) {$\c$};
    \node at ({2.55*cos(\a)},{2.55*sin(\a)}) {$\p$};
  }
  \node[TextGray,align=left,anchor=west] at (4.0,0.45)
    {中心导体孤立且初始中性\\$\sum C_i(\varphi_0-\varphi_i)=0$};
\end{tikzpicture}

{\small\color{TextGray}图 32.16\quad 三电容星形连接的浮置中心节点}
\end{center}}

\newcommand{\atmosphereCapacitorDiagram}{%
\begin{center}
\begin{tikzpicture}[scale=0.83,>=Latex]
  \fill[PaleBlue] (0,0) circle (2.18); \fill[white] (0,0) circle (1.86);
  \fill[SoftGray] (0,0) circle (1.62);
  \draw[curve] (0,0) circle (1.62); \draw[DeepBlue,thick] (0,0) circle (2.18);
  \node at (0,0) {地球};
  \foreach \a in {0,35,...,325}{
    \node[Ink] at ({1.72*cos(\a)},{1.72*sin(\a)}) {$-$};
    \node[DeepBlue] at ({2.05*cos(\a)},{2.05*sin(\a)}) {$+$};
  }
  \draw[<->] (2.35,0)--node[right] {$h$}(2.35,0.56);
  \node[TextGray,align=left,anchor=west] at (4.1,0.7)
    {内半径 $R_\oplus$\\外半径 $R_\oplus+h$\\大气同时提供电容与漏电电阻};
\end{tikzpicture}

{\small\color{TextGray}解图 32.11-A\quad Kelvin 的大气球形电容器模型}
\end{center}}
